\documentclass{amsart}
% For dealing with references we use the comment environment
\usepackage{verbatim}
\newenvironment{reference}{\comment}{\endcomment}
%\newenvironment{reference}{}{}
\newenvironment{slogan}{\comment}{\endcomment}
\newenvironment{history}{\comment}{\endcomment}

% For commutative diagrams we use Xy-pic
\usepackage[all]{xy}

% We use 2cell for 2-commutative diagrams.
\xyoption{2cell}
\UseAllTwocells

% We use multicol for the list of chapters between chapters
\usepackage{multicol}

% This is generally recommended for better output
\usepackage{lmodern}
\usepackage[T1]{fontenc}

% For cross-file-references

% Package for hypertext links:
\usepackage{hyperref}

% For any local file, say "hello.tex" you want to link to please

% Theorem environments.
%
\theoremstyle{plain}
\newtheorem{theorem}[subsection]{Theorem}
\newtheorem{proposition}[subsection]{Proposition}
\newtheorem{lemma}[subsection]{Lemma}

\theoremstyle{definition}
\newtheorem{definition}[subsection]{Definition}
\newtheorem{example}[subsection]{Example}
\newtheorem{exercise}[subsection]{Exercise}
\newtheorem{situation}[subsection]{Situation}

\theoremstyle{remark}
\newtheorem{remark}[subsection]{Remark}
\newtheorem{remarks}[subsection]{Remarks}

\numberwithin{equation}{subsection}

% Macros
%
\def\lim{\mathop{\mathrm{lim}}\nolimits}
\def\colim{\mathop{\mathrm{colim}}\nolimits}
\def\Spec{\mathop{\mathrm{Spec}}}
\def\Hom{\mathop{\mathrm{Hom}}\nolimits}
\def\Ext{\mathop{\mathrm{Ext}}\nolimits}
\def\SheafHom{\mathop{\mathcal{H}\!\mathit{om}}\nolimits}
\def\SheafExt{\mathop{\mathcal{E}\!\mathit{xt}}\nolimits}
\def\Sch{\mathit{Sch}}
\def\Mor{\mathop{\mathrm{Mor}}\nolimits}
\def\Ob{\mathop{\mathrm{Ob}}\nolimits}
\def\Sh{\mathop{\mathit{Sh}}\nolimits}
\def\NL{\mathop{N\!L}\nolimits}
\def\CH{\mathop{\mathrm{CH}}\nolimits}
\def\proetale{{pro\text{-}\acute{e}tale}}
\def\etale{{\acute{e}tale}}
\def\QCoh{\mathit{QCoh}}
\def\Ker{\mathop{\mathrm{Ker}}}
\def\Im{\mathop{\mathrm{Im}}}
\def\Coker{\mathop{\mathrm{Coker}}}
\def\Coim{\mathop{\mathrm{Coim}}}

% Boxtimes
%
\DeclareMathSymbol{\boxtimes}{\mathbin}{AMSa}{"02}

%
% Macros for moduli stacks/spaces
%
\def\QCohstack{\mathcal{QC}\!\mathit{oh}}
\def\Cohstack{\mathcal{C}\!\mathit{oh}}
\def\Spacesstack{\mathcal{S}\!\mathit{paces}}
\def\Quotfunctor{\mathrm{Quot}}
\def\Hilbfunctor{\mathrm{Hilb}}
\def\Curvesstack{\mathcal{C}\!\mathit{urves}}
\def\Polarizedstack{\mathcal{P}\!\mathit{olarized}}
\def\Complexesstack{\mathcal{C}\!\mathit{omplexes}}
% \Pic is the operator that assigns to X its picard group, usage \Pic(X)
% \Picardstack_{X/B} denotes the Picard stack of X over B
% \Picardfunctor_{X/B} denotes the Picard functor of X over B
\def\Pic{\mathop{\mathrm{Pic}}\nolimits}
\def\Picardstack{\mathcal{P}\!\mathit{ic}}
\def\Picardfunctor{\mathrm{Pic}}
\def\Deformationcategory{\mathcal{D}\!\mathit{ef}}

\setlength{\emergencystretch}{3em}
\begin{document}
\title[Stacks Project, Tag 0A9W]{256 --- The dualizing complex of an arbitrary projective bundle}
\maketitle
\noindent Copyright (C) 2005--2025 Johan de Jong. Stacks Project authors. Source: \href{https://stacks.math.columbia.edu/tag/0A9W}{Tag 0A9W}.

\noindent Editorial difficulty note (not author text). Difficulty H2: The proof constructs a candidate comparison map from the canonical top projective-bundle trace. It proves the map invertible by testing against a finite generating collection of twists and verifying the appropriate cohomology vanishing for each twist. The generator-detection argument establishes the relative dualizing object itself rather than applying Serre duality..

\noindent Editorial provenance note (not author text). History: excerpt from revision \texttt{a0444\allowbreak 6e57e\allowbreak c1fbc\allowbreak 252a8\allowbreak 71afc\allowbreak ec775\allowbreak 2fb28\allowbreak 07b14}, duality.tex, lines 3294--3360. Reference presentation was adapted; original-author context and core support blocks were retained or added. The mathematical wording is unchanged. Licensed under GNU FDL 1.2 or later; no invariant sections or cover texts. See ../licenses/COPYING. Context blocks reproduce author definitions and supporting results; they are not additional counted problems.

\begin{lemma}
\label{lemma-twisted-inverse-image}
\begin{reference}
This is almost the same as \href{https://stacks.math.columbia.edu/bibliography}{Bibliography: Neeman-Grothendieck (Example 4.2)}.
\end{reference}
Let $f : X \to Y$ be a morphism between quasi-separated and quasi-compact
schemes. The functor $Rf_* : D_\QCoh(\mathcal{O}_X) \to D_\QCoh(\mathcal{O}_Y)$
has a right adjoint.
\end{lemma}

\begin{lemma}
\label{lemma-upper-shriek-P1}
Let $Y$ be a quasi-compact and quasi-separated scheme.
Let $\mathcal{E}$ be a finite locally
free $\mathcal{O}_Y$-module of rank $n + 1$ with determinant
$\mathcal{L} = \wedge^{n + 1}(\mathcal{E})$.
Let $f : X = \mathbf{P}(\mathcal{E}) \to Y$ be the projection.
Let $a$ be the right adjoint for
$Rf_* : D_\QCoh(\mathcal{O}_X) \to D_\QCoh(\mathcal{O}_Y)$ of
Lemma \ref{lemma-twisted-inverse-image}.
Then there is an isomorphism
$$
c : f^*\mathcal{L}(-n - 1)[n] \longrightarrow a(\mathcal{O}_Y)
$$
In particular, if $\mathcal{E} = \mathcal{O}_Y^{\oplus n + 1}$, then
$X = \mathbf{P}^n_Y$ and we obtain
$a(\mathcal{O}_Y) = \mathcal{O}_X(-n - 1)[n]$.
\end{lemma}

\begin{proof}
In (the proof of) Cohomology of Schemes, Lemma
\href{https://stacks.math.columbia.edu/tag/01XX}{lemma cohomology projective bundle (Tag 01XX)}
we constructed a canonical isomorphism
$$
R^nf_*(f^*\mathcal{L}(-n - 1)) \longrightarrow \mathcal{O}_Y
$$
Moreover, $Rf_*(f^*\mathcal{L}(-n - 1))[n] = R^nf_*(f^*\mathcal{L}(-n - 1))$,
i.e., the other higher direct images are zero. Thus we find an isomorphism
$$
Rf_*(f^*\mathcal{L}(-n - 1)[n]) \longrightarrow \mathcal{O}_Y
$$
This isomorphism determines $c$ as in the statement of the lemma
because $a$ is the right adjoint of $Rf_*$.
By Lemma \href{https://stacks.math.columbia.edu/tag/0A9P}{lemma proper noetherian (Tag 0A9P)} construction of the $a$
is local on the base. In particular, to check that
$c$ is an isomorphism, we may work locally on $Y$.
In other words, we may assume $Y$ is affine and
$\mathcal{E} = \mathcal{O}_Y^{\oplus n + 1}$.
In this case the sheaves $\mathcal{O}_X, \mathcal{O}_X(-1), \ldots,
\mathcal{O}_X(-n)$ generate $D_\QCoh(X)$, see
Derived Categories of Schemes, Lemma \href{https://stacks.math.columbia.edu/tag/0A9V}{lemma generator P1 (Tag 0A9V)}.
Hence it suffices to show that
$c : \mathcal{O}_X(-n - 1)[n] \to a(\mathcal{O}_Y)$
is transformed into an isomorphism under the functors
$$
F_{i, p}(-) = \Hom_{D(\mathcal{O}_X)}(\mathcal{O}_X(i), (-)[p])
$$
for $i \in \{-n, \ldots, 0\}$ and $p \in \mathbf{Z}$.
For $F_{0, p}$ this holds by construction of the arrow $c$!
For $i \in \{-n, \ldots, -1\}$ we have
$$
\Hom_{D(\mathcal{O}_X)}(\mathcal{O}_X(i), \mathcal{O}_X(-n - 1)[n + p]) =
H^p(X, \mathcal{O}_X(-n - 1 - i)) = 0
$$
by the computation of cohomology of projective space
(Cohomology of Schemes, Lemma
\href{https://stacks.math.columbia.edu/tag/01XT}{lemma cohomology projective space over ring (Tag 01XT)})
and we have
$$
\Hom_{D(\mathcal{O}_X)}(\mathcal{O}_X(i), a(\mathcal{O}_Y)[p]) =
\Hom_{D(\mathcal{O}_Y)}(Rf_*\mathcal{O}_X(i), \mathcal{O}_Y[p]) = 0
$$
because $Rf_*\mathcal{O}_X(i) = 0$ by the same lemma.
Hence the source and the target of $F_{i, p}(c)$ vanish
and $F_{i, p}(c)$ is necessarily an isomorphism.
This finishes the proof.
\end{proof}
\end{document}
